\section{Selección de herramientas}
El siguiente paso en el proceso consiste en analizar en profundidad un conjunto reducido de herramientas, seleccionadas a partir de los resultados de la etapa anterior. Si bien el objetivo es analizar herramientas, no resulta viable realizar un análisis en profundidad de la totalidad de las soluciones obtenidas con el filtrado descrito, ya que analizar exhaustivamente cada herramienta implicaría un esfuerzo significativo en términos de tiempo y recursos.

Con el fin de acotar el conjunto de herramientas a analizar y concentrar el estudio en aquellas que presentan un nivel aceptable de funcionamiento, se definió un umbral mínimo de selección sobre el puntaje obtenido en el ranking ponderado descrito previamente. Dicho umbral fue establecido en 0.60, considerando que una herramienta con una puntuación igual o superior a este valor cumple adecuadamente con la mayoría de los criterios esenciales evaluados.

Por el contrario, aquellas herramientas que obtienen un puntaje inferior a este valor presentan deficiencias funcionales relevantes, en la mayoría de los casos mostrando que no cumplen con al menos dos de las funcionalidades estudiadas, por lo que quedan descartadas del proceso.

Es importante destacar que esta etapa no tiene como objetivo determinar de forma definitiva la herramienta más adecuada, sino reducir el conjunto inicial de soluciones a aquellas que cumplen con un nivel mínimo aceptable de funcionalidades, permitiendo así llevar a cabo un análisis más profundo en las herramientas estudiadas.

\paragraph{Herramientas seleccionadas}
La Tabla~\ref{tab:herramientas_hu} resume las herramientas evaluadas, sus principales características y el puntaje obtenido.

\begin{sidewaystable}[p]
    \centering
    \caption{Puntajes obtenidos por herramienta, respectivos al cumplimiento de funcionalidades}
    \label{tab:herramientas_hu}
    \scriptsize
    \renewcommand{\arraystretch}{1.3}

    \begin{tabularx}{\textheight}{|
        X|
        >{\centering\arraybackslash}X|
        >{\centering\arraybackslash}X|
        >{\centering\arraybackslash}X|
        >{\centering\arraybackslash}X|
        >{\centering\arraybackslash}X|
        >{\centering\arraybackslash}X|
        c|}
        \hline
        \textbf{Herramienta} & 
        \textbf{Generación Automática de HU} & 
        \textbf{Funcionamiento automático} &
        \textbf{Criterios de Aceptación} & 
        \textbf{Detección de Inconsistencias} & 
        \textbf{Validación / Refinamiento} & 
        \textbf{Trazabilidad Req/HU} & 
        \textbf{Puntaje} \\ \hline
        ClickUp AI                 & Sí & No & Sí & Sí & Sí & Sí & \textbf{0.85} \\ \hline
        Modern Requirements4DevOps & Sí & No & Sí & Sí & Sí & Sí & \textbf{0.85} \\ \hline
        POPal (Jira)               & Sí & Sí & Sí & No & Sí & No & \textbf{0.65} \\ \hline
        StoryMachine               & Sí & Sí & Sí & No & Sí & No & \textbf{0.65} \\ \hline
        Miro                       & Sí & No & Sí & Sí & No & Sí & \textbf{0.60} \\ \hline
        StoriesOnBoard             & Sí & Sí & No & No & Sí & No & 0.40 \\ \hline
        Specifai                   & Sí & Sí & No & No & Sí & No & 0.40 \\ \hline
        QAssert                    & Sí & Sí & Sí & No & No & No & 0.40 \\ \hline
        AI Requirements Copilot    & Sí & Sí & Sí & No & No & No & 0.40 \\ \hline
        AutoDev Editor             & Sí & No & Sí & No & No & No & 0.25 \\ \hline
        StoryCraft                 & No & Sí & No & No & No & No & 0.00 \\ \hline
        Agile-bot                  & No & Sí & No & No & No & Sí & 0.00 \\ \hline
        StoryBot                   & No & No & No & No & No & No & 0.00 \\ \hline
        Storygenie                 & Sí & No & No & No & No & No & 0.00 \\ \hline
        AutoStory (Jira)           & No & No & No & No & No & No & 0.00 \\ \hline

    \end{tabularx}

\end{sidewaystable}

Las herramientas se presentan ordenadas de forma descendente según el puntaje obtenido, destacándose en negrita aquellas que superan el umbral mínimo definido para la siguiente etapa del análisis.

\begin{comment}

    \paragraph{Herramientas seleccionadas}
La siguiente tabla resume las herramientas evaluadas, sus principales características y el puntaje obtenido, indicando de esta forma aquellas que resultaron seleccionadas para la siguiente etapa del análisis.
\begin{landscape}
\begin{table}[p]
    \centering
    \caption{Evaluación de herramientas para la generación de historias de usuario}
    \label{tab:herramientas_hu}
    \scriptsize
    \renewcommand{\arraystretch}{1.3}

    \adjustbox{max width=\linewidth, max totalheight=\textheight, keepaspectratio}{%
    \begin{tabularx}{\linewidth}{|
        X|
        >{\centering\arraybackslash}X|
        >{\centering\arraybackslash}X|
        >{\centering\arraybackslash}X|
        >{\centering\arraybackslash}X|
        >{\centering\arraybackslash}X|
        >{\centering\arraybackslash}X|
        c|}
        \hline
        \textbf{Herramienta} & 
        \textbf{Generación Automática de HU} & 
        \textbf{Funcionamiento automático} &
        \textbf{Criterios de Aceptación} & 
        \textbf{Detección de Inconsistencias} & 
        \textbf{Validación / Refinamiento} & 
        \textbf{Trazabilidad Req/HU} & 
        \textbf{Puntaje} \\ \hline
        ClickUp AI                 & Sí & No & Sí & Sí & Sí & Sí & \textbf{0.85} \\ \hline
        POPal (Jira)               & Sí & Sí & Sí & No & Sí & No & \textbf{0.65} \\ \hline
        StoryCraft                 & No & Sí & No & No & No & No & 0.00 \\ \hline
        Modern Requirements4DevOps & Sí & No & Sí & Sí & Sí & Sí & \textbf{0.85} \\ \hline
        StoriesOnBoard             & Sí & Sí & No & No & Sí & No & 0.40 \\ \hline
        Miro                       & Sí & No & Sí & Sí & No & Sí & \textbf{0.60} \\ \hline
        StoryMachine               & Sí & Sí & Sí & No & Sí & No & \textbf{0.65} \\ \hline
        Agile-bot                  & No & Sí & No & No & No & Sí & 0.00 \\ \hline
        AutoDev Editor             & Sí & No & Sí & No & No & No & 0.25 \\ \hline
        Specifai                   & Sí & Sí & No & No & Sí & No & 0.40 \\ \hline
        StoryBot                   & No & No & No & No & No & No & 0.00 \\ \hline
        QAssert                    & Sí & Sí & Sí & No & No & No & 0.40 \\ \hline
        AI Requirements Copilot    & Sí & Sí & Sí & No & No & No & 0.40 \\ \hline
        Storygenie                 & Sí & No & No & No & No & No & 0.00 \\ \hline
        AutoStory (Jira)           & No & No & No & No & No & No & 0.00 \\ \hline
    \end{tabularx}}
\end{table}
\end{landscape}


\end{comment}

